\documentclass[10pt,a4paper]{amsart}
\usepackage[margin=19mm]{geometry}
\usepackage{amsmath,amssymb,mathtools}
\usepackage[scr=rsfs,cal=euler]{mathalfa}
\usepackage{xfrac}
\usepackage[unicode,hidelinks,bookmarks=false]{hyperref}
\newcommand{\ie}{i.e.\ }
\newcommand{\cf}{cf.\ }
\newcommand{\resp}{respectively\ }
\newcommand{\Subsection}{\subsection}
\providecommand{\nobreakdash}{\nobreak\hbox{-}}
\setlength{\emergencystretch}{2em}
\multlinegap0pt
\usepackage{bm}
\usepackage{bbm}
\usepackage{dsfont}
\usepackage{xspace}
\usepackage{cancel}
\usepackage{mathtools}
\usepackage{amsthm}
\makeatletter
\@ifundefined{Theorem}{\newtheorem{Theorem}{Theorem}}{}
\@ifundefined{Lemma}{\newtheorem{Lemma}[Theorem]{Lemma}}{}
\@ifundefined{Observation}{\newtheorem{Observation}[Theorem]{Observation}}{}
\makeatother
\title[Ordinals and recursively defined functions on the reals]{Ordinals and recursively defined functions on the reals}
\author{Gabriel Nivasch, Lior Shiboli}
\date{Source: arXiv:2311.17210v5 (CC BY 4.0)}
\begin{document}
\maketitle

\noindent\emph{Source and licence.} Ordinals and recursively defined functions on the reals. Version arXiv:2311.17210v5 (CC BY 4.0), distributed under the Creative Commons Attribution 4.0 International licence, as recorded in the arXiv licence metadata for this exact version and shown on the abstract page https://arxiv.org/abs/2311.17210v5. Full rights notice: licenses/arxiv-2311.17210-CC-BY-4.0.txt.\par\smallskip

\noindent\textit{Editorial scope.} Counted unit: the Lemma whose source label is \texttt{lem\_num\_intervals} in the article (source0.tex lines 341--343, source label \texttt{lem\_num\_intervals}), delivered together with the article's own complete original proof of it (source0.tex lines 345--378, one \texttt{proof} environment of 34 lines). No mathematics is written, repaired or re-derived: statement and proof are the article's own text line for line. Required context retained uncounted: none. Every definition, display and label of those lines is retained unchanged. Omitted from this package and not redistributed: 1--340, 344--344, 379--682. Nothing inside a retained excerpt is omitted or altered: every retained body is delivered exactly as the article's lines are written, so no omission receipt of any kind accompanies this package. Citations: the retained excerpts carry no citation command at all, so no bibliography entry of the article is transcribed and no reference is redistributed. Reference adaptation: none was needed and none was made; every reference inside the delivered unit resolves inside it, so no unresolved reference or citation is produced. Printed slips kept literal: no printed word, symbol, hypothesis or number of the counted statement or of its proof was altered. Layout and class: the article's own class is \texttt{article} with packages including inputenc, geometry, amsfonts, amsmath, amssymb, amsthm, xcolor, graphicx, algpseudocode; the wrapper is the lane's pinned \texttt{amsart} editorial wrapper at 10pt on A4, which restates the theorem-like environments the retained text needs and adds the article's own macro definitions that the retained text needs (none), with extra wrapper packages (bm, bbm, dsfont, xspace, cancel, mathtools). The delivered numbering is the unit's own. Assets: no retained span contains a figure environment, a graphics-inclusion command, a PDF-page inclusion, a table or a TikZ picture; no image file is copied into this package, the package declares no asset dependency and no image file is redistributed with it. Licence: version arXiv:2311.17210v5 of this article is distributed under the Creative Commons Attribution 4.0 International licence, as recorded in the arXiv licence metadata for this exact version and shown on the abstract page https://arxiv.org/abs/2311.17210v5. A full rights notice is delivered at licenses/arxiv-2311.17210-CC-BY-4.0.txt. Characters: every retained excerpt is pure ASCII, so the delivered engine, XeLaTeX with its default Latin Modern text font, reports no missing character.
\begin{Lemma}\label{lem_num_intervals}
    The ordinal number of intervals $I_\alpha$ into which $D$ is partitioned is at most $\omega\cdot(o(f|_D)+1)$.
\end{Lemma}

\begin{proof}
Recall that $f$ is positive for all $x\in D$. Call $x\in D$ a \emph{near-root} of $f$ if there exists an infinite increasing sequence $y_1,y_2,y_3,\ldots\in D$ such that $\lim_{n\to\infty} y_n = x$ and $\lim_{n\to\infty} f(y_n) = 0$. Let $L_f\subset \mathbb R$ be the set of near-roots of $f$. The set $L_f$ is well-ordered in $\mathbb R$, since from an infinite decreasing sequence of near-roots $x_1>x_2>x_3>\cdots$ we could construct an infinite $f$-bad sequence $z_1>z_2>z_3>\cdots$ by letting $z_1=x_1$, and for each $i=2,3,\ldots$ choosing $z_i\in(x_i,x_{i-1})$ such that $f(z_i)<f(z_{i-1})$. Hence, the set $L_f$ has an ordinal type $o(L_f)$.

Furthermore, the set $L_f$ is closed: If $x_1<x_2<x_3<\cdots$ are near-roots of $f$ that converge to $x$, then $x$ is a near-root of $f$ as well. Hence, the elements of $L_f$ can be classified into \emph{limit} and \emph{non-limit} near-roots. 

Let $\alpha\mapsto x_\alpha$, for $\alpha<o(L_f)$, be the order-preserving enumeration function of $L_f$. Note that $\alpha$ is a limit ordinal if and only if $x_\alpha$ is a limit near-root of $f$.

We will now show that $o(f|_D)\ge o(L_f)$.
For this, we prove by ordinal induction on $\alpha$ that the following is true for every $\alpha<o(L_f)$: For every $\beta<\alpha$ and every $\varepsilon>0$ there exists $y$ such that:
\begin{itemize}
    \item $x_\alpha-\varepsilon<y<x_\alpha$,
    \item $f(y)<\varepsilon$, and
    \item $o_f(y)\ge \beta$.
\end{itemize}
Indeed, given $\alpha$, given $\beta<\alpha$, and given $\varepsilon>0$, we can choose $y$ close enough to $x_\alpha$ such that $y\ge x_\beta$, $y>x_\alpha-\varepsilon$, and $f(y)<\varepsilon$. Then, applying the induction assumption on $\beta$, for every $\gamma<\beta$ we can find $z<x_\beta$ such that $f(z)<f(y)$ and $o_f(z)\ge \gamma$. It follows that $o_f(y)\ge \beta$ as desired.

Hence, we indeed have $o(f|_D)\ge o(L_f)$, as claimed.

Next, we determine the relation between the intervals $I_\alpha$ and the elements of $L_f$:


\begin{Observation}\label{obs_w-many}
    Let $z\in D$. Then there exists an infinite sequence of $\omega$ consecutive intervals $I_{\alpha}, I_{\alpha+1},\allowbreak I_{\alpha+2},\ldots$ that converge to $z$ if and only if $z$ is a non-limit near-root of $f$.
\end{Observation}

\begin{proof}
Suppose first that $z$ is not a near-root of $f$. Then there exists an $\varepsilon\in(0,x_{\mathrm{hi}}-z)$ such that $f(x)\ge\varepsilon$ for all $x\in [z-\varepsilon,z+\varepsilon]$ (where the part $x\in[z-\varepsilon, z]$ follows from the fact that $z$ is not a near-root, and the part $x\in[z, z+\varepsilon]$ follows from the fact that $f$ is ordinal decreasing). Therefore, an interval $I_\alpha$ whose left endpoint is in $(z-\varepsilon,z)$ must contain $z$ in its interior. Hence, there are not $\omega$-many intervals converging to $z$.

Now suppose $z$ is a near-root of $f$. Then no interval $I_\alpha$ whose left endpoint is left of $z$ can contain $z$ in its interior. If $z$ is a non-limit near-root of $f$, then there exists an $\varepsilon>0$ such that $(z-\varepsilon,z)$ contains no near-roots of $f$. Hence, for every $\varepsilon'<\varepsilon$, the interval $(z-\varepsilon,z-\varepsilon')$ contains only finitely many intervals $I_\alpha$. And therefore, there exist $\omega$-many consecutive intervals $I_\alpha$ converging to $z$. If, on the other hand, the near-root $z$ is itself a limit of near-roots of $f$, then some left-neighborhood of $z$ contains at least $\omega^2$-many intervals $I_\alpha$.
\end{proof}

Observation \ref{obs_w-many} implies that there is a one-to-one correspondence between non-limit elements of $L_f$ and sequences of $\omega$-many consecutive intervals $I_\alpha$, except for a possible final sequence of at most $\omega$-many intervals after the last element of $L_f$.
Lemma \ref{lem_num_intervals} follows.
\end{proof}

\end{document}
